\documentclass[10pt]{article}
\usepackage[margin=0.72in]{geometry}
\usepackage{lmodern}
\usepackage{amsmath,booktabs,array,enumitem,xcolor,hyperref}
\hypersetup{colorlinks=true,linkcolor=blue,urlcolor=blue}
\setlist[enumerate,itemize]{nosep,leftmargin=*}
\setlist[description]{nosep}
\setlength{\parskip}{4pt}
\setlength{\parindent}{0pt}
\newcommand{\ui}[1]{\textsf{\textbf{#1}}}
\newcommand{\code}[1]{\texttt{#1}}
\title{Multi-Tier Digital Twin v2\\Workflow and Calculation Guide}
\author{Rao Yenamandra --- \texttt{raosy@digitaltwinsim.com}\\
Mubanga Nsofu --- \texttt{mubanga.nsofu@vodacom.co.za}}
\date{Version 2 --- September 2026}

\begin{document}
\maketitle

\section{Purpose and evidence boundary}
The console is a screening-grade research model of a heterogeneous radio network containing low-band and mid-band macro cells, an upper-mid-band 6G layer, mmWave, indoor Wi-Fi~7, and a LEO satellite overlay. It compares static, demand-following, rule-based, PPO, and Double-DQN control and can search selected mobility and steering thresholds.

The public GitHub Pages console contains a \textbf{self-contained JavaScript v2 screening engine}. Input changes recalculate capacity, run an aggregate synthetic policy comparison, train lightweight linear PPO and Double-DQN approximations, and search selected knobs locally. This browser model is intentionally simpler than the Python twin and does not claim numerical parity. A separately regenerated five-seed Python dataset is embedded as a reference-validation record. Results use synthetic sessions and are advisory, not carrier-performance or standards-conformance evidence.

\textbf{Important mIoT boundary:} v2 represents offered mIoT device load and scheduled data resources, but it does not model PRACH preamble contention, collisions, retries, access-class barring, Random Access Response capacity, or access failures. Its mIoT result must not be presented as PRACH-storm performance.

\section{End-to-end workflow}
\begin{enumerate}
  \item \textbf{Choose a scenario.} This selects area, session mix, mobility, weather, and available tiers.
  \item \textbf{Build the deployment.} Cell positions, spectrum, PRBs, propagation parameters, MIMO gain, and steering constraints are resolved.
  \item \textbf{Generate sessions.} Each synthetic session receives a service type (eMBB, URLLC, or mIoT), position, indoor state, velocity, and offered bit rate.
  \item \textbf{Run the fast loop.} Every tick, sessions move; link levels, shadow fading, interference, A3/A5 conditions, and handovers are updated.
  \item \textbf{Run the slow loop.} At each control interval, the selected policy chooses one slice-allocation template per tier and one network-wide steering posture.
  \item \textbf{Allocate service and score it.} The twin calculates served traffic, satisfaction, utilization, fairness, energy, handovers, continuity, and reward.
  \item \textbf{Compare policies.} All policies use the same scenario and evaluation seed. A learned policy is recommended only if its gain is material and continuity does not regress.
  \item \textbf{Optionally optimize thresholds.} A one-pass coordinate search tests a finite grid of knob values for the selected objective.
  \item \textbf{Validate across seeds.} v2 repeats PPO and Double-DQN training/evaluation for five seeds and reports means and 95\% Student-$t$ intervals. This measures seed sensitivity, not convergence.
\end{enumerate}

\section{Buttons: exactly what they do}
\begin{tabular}{>{\raggedright\arraybackslash}p{0.25\textwidth}>{\raggedright\arraybackslash}p{0.68\textwidth}}
\toprule
Control & Operation\\
\midrule
\ui{Preview capacity \& coverage} & Constructs the edited deployment and recalculates link-budget coverage and theoretical/loaded capacity without training PPO or DQN.\\
\ui{Run comparison} & Trains the browser PPO and Double-DQN approximations, evaluates them and three deterministic controllers on the entered seeds, and applies the decision guard against the best deterministic controller.\\
\ui{Run optimizer} & Performs a local one-pass coordinate search over the knob grid for the selected objective.\\
\ui{Run full analysis} & Executes preview, comparison, five-seed sensitivity, and selected-objective optimization in sequence.\\
\ui{Reset to scenario defaults} & Live version only: restores the tier table. It does not reset trained weights outside the next requested run.\\
\bottomrule
\end{tabular}

\section{Scenario controls}
\begin{tabular}{>{\raggedright\arraybackslash}p{0.25\textwidth}>{\raggedright\arraybackslash}p{0.47\textwidth}>{\raggedright\arraybackslash}p{0.20\textwidth}}
\toprule
Knob & Meaning and effect & Unit / range\\
\midrule
\ui{Scenario} & Selects the built-in deployment and traffic profile: urban dense, highway mobility, indoor enterprise, rural coverage hole, heavy rain, or cellular only. & Categorical\\
\ui{Sessions} & Number of simultaneous synthetic UE sessions. More sessions raise offered load and possible congestion. & Count $\geq 1$\\
\ui{Indoor fraction} & Probability that a session is indoors; affects penetration loss and eligibility/preference for Wi-Fi. & 0--1\\
\ui{Control intervals} & Number of slow policy decisions in one episode. & Positive integer\\
\ui{Ticks per interval} & Number of mobility and handover updates between policy decisions. & Positive integer\\
\ui{Tick (s)} & Duration of one fast-loop update. A control interval lasts tick duration times ticks per interval. & Seconds\\
\ui{Rain rate} & Drives frequency-dependent rain attenuation, especially on high-frequency and satellite paths. & mm/h $\geq0$\\
\ui{Doppler pre-compensation} & Fraction of predicted NTN Doppler removed before residual inter-carrier-interference loss is calculated. & 0--1\\
\ui{Training episodes} & Number of complete episodes used separately for PPO and Double DQN. More episodes increase experience but do not guarantee convergence. & Positive integer\\
\ui{Seed} & Controls reproducible session generation, fading, training exploration, and the held-out evaluation seed offset. & Integer\\
\bottomrule
\end{tabular}

These fields are editable in the v2 browser console. A changed value affects the next local calculation but does not overwrite the embedded Python reference dataset.

\section{Per-tier radio controls}
\begin{description}
  \item[Cells] Number of deployed cells in the tier. Zero removes the tier. It changes aggregate capacity, geometry, interference, and coverage opportunities.
  \item[Morphology] Selects the TR~38.901-shaped UMa, UMi, RMa, or InH path-loss function. LEO uses the NTN link-budget path.
  \item[Bandwidth] Channel bandwidth per component carrier in MHz.
  \item[Carriers] Number of aggregated component carriers. Aggregate bandwidth is bandwidth times carriers.
  \item[MIMO layers] Spatial-layer count used in the screening capacity multiplier.
  \item[Mode] SU or MU. MU applies an additional efficiency factor in the simplified capacity model.
\end{description}

The corrected NR grids explicitly align SCS and PRB width: macro-low uses 15 kHz/0.18 MHz, macro-mid and NTN use 30 kHz/0.36 MHz, upper-mid-band uses 60 kHz/0.72 MHz, and mmWave uses 120 kHz/1.44 MHz. Wi-Fi uses this generic field as a 2 MHz scheduling-resource width rather than an NR PRB.

\section{Core radio and capacity calculations}
For a derived grid, the bandwidth represented by one PRB is
\[
B_{\mathrm{PRB,MHz}}=\frac{12\,\Delta f_{\mathrm{kHz}}}{1000},\qquad
N_{\mathrm{PRB}}=\left\lfloor\frac{0.95 B_{\mathrm{MHz}}N_{\mathrm{CC}}}{B_{\mathrm{PRB,MHz}}}\right\rfloor .
\]
Explicit standard/reference grids override this approximation. The screening MIMO multiplier is
\[
G_{\mathrm{MIMO}}=\max(1,L)\times0.75\times
\begin{cases}1.15,&\text{MU mode}\cr1,&\text{SU mode.}\end{cases}
\]
The displayed theoretical peak per cell is
\[
C_{\mathrm{peak}}=N_{\mathrm{PRB}}B_{\mathrm{PRB,MHz}}\eta_{\max}G_{\mathrm{MIMO}}\quad\text{Mbit/s}.
\]
This is an upper bound, not measured throughput. Achieved service uses the simulated SINR, interference coupling, attachment, mobility derating, demand, and policy allocation.

The maximum allowable path loss is formed from transmit power, antenna gains, thermal noise, receiver noise figure, required SINR, and margin:
\[
\mathrm{MAPL}=P_{\mathrm{tx}}+G_{\mathrm{tx}}+G_{\mathrm{rx}}
-\bigl[-174+10\log_{10}(B_{\mathrm{PRB,Hz}})+NF\bigr]
-\mathrm{SINR}_{\mathrm{req}}-M.
\]
The terrestrial radius is the distance at which the morphology-specific path loss reaches MAPL. NTN reports orbital slant range. The displayed radius is a screening link-budget result and should not be confused with an engineered inter-site distance.

For slice $s$, satisfaction is clipped to the interval 0--1:
\[
q_s=\min\!\left(1,\frac{\text{served traffic}_s}{\max(\text{offered traffic}_s,\epsilon)}\right).
\]
Jain fairness is $J=(\sum_s q_s)^2/(3\sum_s q_s^2+\epsilon)$.

\section{Mobility, handover, and steering}
The fast loop evaluates biased serving and neighbour measurements. A3 requires a neighbour to exceed the serving cell by the configured offset and hysteresis for the time-to-trigger. A5 combines a weak-serving condition with a sufficiently strong target. A completed attempt may succeed, fail, or later count as ping-pong according to the model's parametric execution rules.

Key mobility/steering knobs used by the optimizer are:
\begin{itemize}
  \item A3 offset: 1, 2, 3, 4.5, or 6 dB;
  \item A3 hysteresis: 0.5, 1, 2, or 3 dB;
  \item A3 time-to-trigger: 40, 100, 160, 320, or 640 ms;
  \item A5 serving and neighbour absolute RSRP thresholds (editable in dBm);
  \item cell individual offset: $-6$, 0, 4, 8, 12, or 16 dB;
  \item NTN speed threshold: 60, 90, 120, 160, or 220 km/h;
  \item NTN coverage-hole threshold: $-126$, $-120$, $-114$, or $-108$ dBm;
  \item Wi-Fi indoor bias: 2, 6, 10, or 14 dB;
  \item minimum RSRP: $-135$, $-130$, $-125$, or $-120$ dBm.
\end{itemize}

\section{Policies, actions, and reward}
One centralized agent observes the complete network. It emits one of nine slice-mix templates for every active tier plus one of four global steering postures. With six tiers this branching representation corresponds to $9^6\times4=2{,}125{,}764$ possible joint actions without constructing one enormous flat output table.

The five comparisons are non-adaptive fixed allocation, demand-following allocation, a transparent adaptive rule, PPO, and Double DQN. The non-adaptive controller repeats the same 55/25/20 eMBB/URLLC/mIoT allocation at every control interval even though radio conditions continue to change. PPO is an on-policy clipped policy-gradient method; Double DQN learns branch-wise action values using replay and a target network.

At each control interval the scalar reward is
\[
r=\mathbf{w}^{T}\mathbf{q}+0.06U+0.04J-0.12V
-c_EE-c_HH-c_C(1-C),
\]
where $\mathbf q$ is slice satisfaction, $U$ utilization, $J$ fairness, $V$ weighted SLA violation, $E$ normalized energy, $H$ handover churn, and $C$ session continuity. Reward is a model preference function, not a physical KPI.

The browser recommendation guard first selects the better learned controller by reward and compares it with the best of the three deterministic controllers. It rejects the learned controller if continuity falls more than 0.5 percentage points below that comparator, if aggregate gain is negative, or if gain is below the 2\% materiality threshold. A PASS therefore means ``preferred within this synthetic run,'' not universally superior. The Python reference retains its separately documented rule-based comparison rule.

\section{Optimizer objectives}
The optimizer tests one knob at a time over its finite grid, retains an improvement, and moves to the next knob. One pass is not a global optimum.
\begin{itemize}
  \item \ui{Handover thresholds}: rewards continuity and handover success; penalizes ping-pong and failures.
  \item \ui{Throughput (eMBB)}: maximizes total served Gbit/s plus eMBB satisfaction, subject to continuity and ping-pong penalties.
  \item \ui{Latency (URLLC)}: minimizes mean modeled access latency while retaining URLLC satisfaction and continuity.
  \item \ui{Coverage (mIoT / edge)}: maximizes covered offered traffic and continuity. It is radio coverage, not PRACH accessibility.
\end{itemize}
If the feasibility panel says coverage-limited or capacity-limited, threshold tuning is not the correct remedy; sites, spectrum, low band, or NTN may be required.

\section{How to read the v2 outputs}
\begin{description}
  \item[Feasibility] Separates controller-tunable cases from missing coverage/capacity.
  \item[Decision status] Applies the material-gain and continuity guard against the rule-based controller.
  \item[KPI comparison] Gives exact reward, per-slice satisfaction, handover, continuity, latency, energy, and gain values.
  \item[Visual comparison] Shows selected policy bars; use the KPI table for exact numbers and all five policies.
  \item[Five-seed sensitivity] The browser run reports means and 95\% Student-$t$ intervals over the editable seed list, plus direct PPO-versus-DQN reward wins. The embedded Python record uses seeds 3080--3084. Five seeds are useful for detecting instability but are not enough to prove convergence.
  \item[Optimization] Shows baseline and selected knob values, objective change, number of evaluations, and search time.
  \item[Provenance] Records model version, UTC generation time, generator, and a configuration hash. Identical labels do not imply identical data unless the hash also matches.
\end{description}

\section{Minimum review checklist}
Before citing a result, record the scenario, seed set, configuration hash, controller, training episodes, evaluation rule, and confidence interval. Confirm that the recommended policy does not reduce continuity, distinguish theoretical peak from achieved service, and state that PRACH is absent. For stronger evidence, increase training episodes, plot learning curves, use more independent seeds, reserve unseen test scenarios, test sensitivity to reward weights, and compare against carrier or calibrated traces.

\end{document}
